\section{Conclusions} \label{sec:conclusion}

The oil industry is of utmost importance to the global economy and requires constant optimization mechanisms in its various production phases. In particular, oil extraction in southern Argentina requires periodic reconfiguration of various batteries to reduce operating costs, which in many cases is performed manually by an operator. In this ongoing work, an integer linear programming model is presented to schedule the adjustment of the battery regime by modifying the production regimes of the wells along with the itinerary the operator must follow. The computational experiment shows the validity of the model and the feasibility of using it with open-source solvers.Next steps in this line of work will focus on evaluating the complexity of the model and studying alternatives such as heuristics or meta-heuristics for its resolution. We will also seek to scale the application to a set of batteries with a larger number of wells and, on the other hand, consider technical details regarding pumping systems, such as those that impose restrictions on the feasible values of the working regimes.